\documentclass{article}
\usepackage[utf8]{inputenc}
\usepackage{enumitem}
\usepackage{dsfont}

\usepackage[acronym,nomain,nonumberlist]{glossaries-extra}

\makenoidxglossaries


\newglossaryentry{dynamicaperture}{
    name={dynamic aperture},
    description={Region of stable transverse phase space}
}

\title{Placement progress and notes}
\author{Madeleine Watson}

\begin{document}

\maketitle

\section{Monthly Progress}
\subsection*{Month 1: 23/02/2026 - 20/03/2026}
\subsubsection*{Overview}

Spent the first week familiarising myself with beam optics concepts and Xsuite functions. Then began understanding and editing a code for a compact ring with 90$^\circ$ phase advance.

\subsubsection*{Work completed}

\begin{itemize}
\item Initially read papers to understand how the ring functions and what processes go into it.
\item Made notes on initial ring design code and started compiling work into a GitHub repository.
\item Added sextupoles to the ring design and performed chromaticity correction.
\item Corrected beta-functions to create uniform Twiss plot.
\item Corrected working point and found two potential candidates.
\end{itemize}

\subsubsection*{Problems encountered}

\textbf{Breaking the ring:} \texttt{Invalid n1}, \texttt{Invalid n2}, and \texttt{Invalid n3}errors arose when trying to add sextupoles. This scenario halted progress the most.
\begin{itemize}
    \item \textbf{Solution:} There was no set solution to this error but there were some common causes: changing the geometry of the ring, not matching/rematching the quadrupole strengths. The `breakage' could be identified by computing a Twiss with specified conditions.
\end{itemize}

\textbf{Sextupole positioning:} This would cause ring breakage. It took a week or so to identify sextupole positions that worked.
\begin{itemize}
    \item \textbf{Solution:} I tried creating a function that identified the best positions for sextupoles, but the best solution was trial and error and knowledge of ring properties.
\end{itemize}

\textbf{Large beta-functions:} Changing parameters resulted in increased beta-functions
\begin{itemize}
    \item \textbf{Solution:} Consistently running matching functions and identifying best positions and working points.
\end{itemize}

\subsubsection*{Results and progress}
\begin{itemize}
    \item The initial compact ring now has sextupoles inserted: there are currently 16 in each arc (8 focusing, 8 defocusing).
    \item Strengths and beta-functions have been corrected
    \item Two potential working points have been identified and plotted on resonance graphs to show suitability
    \item Using these working points resulted in reduced beta-functions, more uniform Twiss plots, and slightly reduced sextupole strengths.
    \item Another configuration with a decreased triplet drift length is being explored and shows signs of reduced beta-functions
    \item Chromatic functions are plotted and show slightly different, but periodic, results for both working points.
\end{itemize}

\subsubsection*{Next steps}
\begin{itemize}
    \item Further refine suitability of the working point
    \item Track off-momentum particles
    \item Track dynamic aperture and momentum acceptance
    \item Work on different sextupole configurations with different positions, families, and cell lengths in parallel
\end{itemize}

\subsubsection*{Skills developed}
\begin{itemize}
    \item \textbf{Knowledge of beam optics concepts:} Twiss parameters, chromaticity, tune, basics of synchrotron radiation, and dynamic aperture
    \item \textbf{Xsuite:} Building rings and inserting elements, calculating Twiss plots, creating matching functions
\end{itemize}

\subsection*{Month 2: 23/03/2026 - 24/04/2026}
\subsubsection*{Overview}

Completed dynamic aperture and momentum acceptance studies for different configurations. Experimented with slicing, integrators, and fringe fields. Now introducing misalignments and corrections.

\subsubsection*{Work completed}

\begin{itemize}
\item Plotted DA/MA for 7 different sextupole configurations 
\item Optimised configurations to increase DA/MA
\item Changed number of slices in integrators to increase quality of plots
\item Introduced fringe fields for every magnet type and compared with "perfect" results
\item Performed frequency map analysis
\item Introduced misalignments, inserted BPMs and correctors, and performed orbit correction
\end{itemize}

\subsubsection*{Problems encountered}

\textbf{Reducing non-linearities:} Chromatic correction is working for $dqx$ and $dqy$ but often plotting tune against momentum deviation showed that $ddqx$ and $ddqy$ had not been properly corrected. This contributed to smaller DA/MA. Performing chromaticity correction to get these values to 0 often failed as it was too much of a jump.
\begin{itemize}
    \item \textbf{Solution:} The chromaticity correction tries to match to 80\% of the value of $ddqx$ and $ddqy$. If this is successful, it iterates through 80\% of the resulting values until the code fails. It takes the strengths from the iteration with the smallest second order chromaticity.
\end{itemize}


\textbf{Code organisation:} Many scripts and results in different places
\begin{itemize}
    \item \textbf{Solution:} Everything has been modularised so it is easier to see where errors arise and can edit things more efficiently. This will also help with later lattice designs.
\end{itemize}

\subsubsection*{Results and progress}
\begin{itemize}
    \item Increasing the number of slices for the quadrupoles created smoother curves in the DA/MA plots, but there were no differences when increasing the slices for bends and quadrupoles.
    \item Fringe fields reduced the less stable configurations' DA/MA considerably, but not as much for stable configurations.
    \item Added another focusing sextupole pair to the configuration with two sextupole pairs at 180$^\circ$ phase advance, which dramatically improved momentum acceptance.
    \item Frequency map analysis calculated and shows stable diffusion
    \item BPMs and correctors are currently placed around every quadrupole (horizontal and vertical) and the MICADO algorithm is used for corrections
    \item With corrections, the orbits are currently oscillating at less than 1mm
\end{itemize}

\subsubsection*{Next steps}
\begin{itemize}
    \item Studies on spin tracking and polarization build-up time
    \item Looking at integration with the rest of the injector complex using current macroparticle distribution
    \item Assess injection efficiency
\end{itemize}

\subsubsection*{Skills developed}
\begin{itemize}
    \item \textbf{Dynamic aperture and momentum acceptance:} how it is calculated and the effects of adjusting integrators and slices
    \item \textbf{Knowledge of beam optics concepts:} fringe fields, FMA, BPMs, correctors, synchrotron radiation damping
\end{itemize}

\glsaddall

\printglossary[title=Glossary]
\end{document}